% r01 The rotation part: lattices, walled tori and the route
\section{The rotation part: lattices, walled tori and the route}\label{sec:rot}

This and the following six sections form the rotation part of the paper. Its object is King's printed state on curvature sources, $z\mapsto\ZKinf(\sum_jz_jf_j,0)$, the entire function of the coefficients of real $2$-forms $f_j\in C^\infty_c$ whose derivatives at $z=0$ are the curvature Schwinger functions (\Cref{hyp:s03:EC} and the curvature-sector section); its result, \Cref{rot:thm:main}, is that this function is invariant under every Euclidean motion.

Rotation invariance is not automatic, since all lattices in King's construction have the same orientation. Theorem~2.4 of~\cite{K2} states that on the tori $\DJ$ the lattices $\eps_k\Z^3$ and $R(\eps_k\Z^3)$, with $R$ the rotation by $\thz=\arctan(3/4)$ in a coordinate plane, have the same continuum limit; since $\thz/\pi$ is irrational these rotations generate a dense subgroup of $SO(3)$. As printed the statement is false in two ways (\Cref{rot:prop:scalar}, \Cref{rot:rem:log}): the source it pairs with the plaquettes is one scalar function on both lattices, which couples the field to two different $2$-forms, and the Wick constant it uses misses an order-$e^2$ logarithm. We therefore call the printed statement \emph{King's torus conjecture}, and we use instead the corrected statement, for $2$-form sources sampled in each lattice's own frame and for the two lattices at the same spacing; it is the torus comparison \Kz\ of \Cref{rot:hyp:K}, proved in \cref{ssec:rot:Kproof}. Moreover, even granting the torus comparison, two steps remain between it and rotation invariance of King's state: the infinite-volume limit along tori, where correlation inequalities give no monotonicity, and its identification with the limit along cubes, which King writes can be expected from a cluster expansion~\cite[(2.24)]{K2}. The ultraviolet analysis of~\cite{K1,K2} needs boundaries that are lattice hyperplanes~\cite[p.~326]{K2}, and no bounded region has that property for both lattices (\Cref{rot:prop:forced}). Where to place the two-lattice comparison is therefore the whole question. The rotation part places it on periodic tori carrying a smooth mass wall, and proves $O(3)$ invariance of King's printed state in the curvature sector from three inputs, each stated in one line and proved where it is stated: the existence theorem \EC\ of \Cref{hyp:s03:EC}, the torus comparison \Kz\ (\Cref{rot:hyp:K}) and the heavy-shell estimates \HS, \DSh\ (\Cref{rot:hyp:HS,rot:hyp:DS}).

\subsection{The torus route}

On each lattice separately King's cubes are compared with tori carrying a \emph{smooth mass wall}. Fix $R$ in the generating set $\Rgen=\{R_{\mu\nu}(\pm\thz)\}$ of \Cref{rot:def:geom}, a pair $(f,g)\in\cK$ (a real $2$-form $f$ and a function $g\ge0$, both smooth with compact support) and a scheme $a$; at $g=0$ nothing below depends on the scheme. For a walled torus $\fc=(J',w)$, the periodic model on $D_{J'}$ in which the scalar mass squared is raised by $2w$, with $w\ge0$ smooth and vanishing near the supports (\Cref{rot:def:tori}, \Cref{rot:def:walls}), the chain is
\begin{equation}\label{rot:eq:chain}
Z^a_{k,L}(\mathrm i\,R\cdot f,R\cdot g)\;=\;Z^a_{R^{-1}\Lk,\,R^{-1}X_L}(\mathrm if,g)\;\le\;\Tw{a}{\fc,R^{-1}\Lk}(f,g)
\;\xrightarrow[k\to\infty]{}\;\Wv^a_\fc(f,g)\;\xleftarrow[k\to\infty]{}\;\Tw{a}{\fc,\Lk}(f,g).
\end{equation}
Here $Z^a_{k,L}$ is the functional of the comparison cube $X_L=C_L\cap\Lk$ on the axis lattice (\Cref{def:s02:sources}); the first equality is a passive rotation (\Cref{rot:lem:walledalg}(b)); the inequality is a correlation inequality on the rotated lattice $R^{-1}\Lk$ (\Cref{rot:lem:cutwall}), for which $R^{-1}C_L$ is a cube with faces on lattice hyperplanes; $\Tw{a}{\fc,\Lambda}$ is the walled ratio of \Cref{rot:def:tori}; and the two continuum limits exist and agree by the torus comparison \Kz\ (\Cref{rot:hyp:K}; at $g\ne0$ its scheme form \Ka, \Cref{rot:hyp:Ka}). The comparison cubes on the left are squeezed, as $k\to\infty$ and then $L\to\infty$, to King's printed state by the interlacing of the two cube families (the section on the interlacing). What is missing is an upper bound for walled tori on the axis lattice by King's printed state:
\begin{equation}\label{rot:eq:missing}
\text{for every }\eta>0\text{ there is an admissible }\fc\text{ with }\Wv_\fc(f,0)\le\ZKinf(\mathrm if,0)+\eta .
\end{equation}
This is the closing estimate \TIz\ of \Cref{rot:hyp:TIGHT}. It involves one lattice and one orientation, and the reverse inequality holds in the limit (\Cref{rot:lem:free}), so \TIz\ says that King's printed state is the lower envelope of the walled tori. With it, letting $L\to\infty$ in~\eqref{rot:eq:chain} gives $\ZKinf(\mathrm i\,R\cdot f,0)\le\ZKinf(\mathrm if,0)$, and applying this to $R^{-1}\in\Rgen$ gives equality. The passage from the dense subgroup to $O(3)$ uses the Lipschitz bound of \Cref{cor:s06:TK} at imaginary coefficients and the cubic symmetry of King's printed state (\Cref{rot:thm:curv}), so no regularity of the Schwinger functions beyond that corollary is needed.

The closing estimate \TIz\ does not follow from the correlation inequalities: \Cref{rot:rem:gauss} gives a family of Gaussian functionals with all the other properties used in the proof whose analogue of King's state is not rotation invariant. We therefore isolate the local estimate that implies \TIz. In the walled torus let the sources sit in a box $\Hbox$ and let $U$ be a thick shell around $\Hbox$ inside the heavy plateau, where the wall is at least $1$. Raise the scalar mass squared on the shell further by $2K$, through a second smooth profile, and let $\Psik(n,K)$ be the resulting expectation at lattice spacing $\eps_k$, with $\Psik(n,\infty)$ the expectation with the matter removed on the shell (\Cref{rot:def:twolevel}). The decoupling property \REz\ of \Cref{rot:hyp:RES} says that for one such configuration and one finite $K$
\[
\liminf_{k\to\infty}\big[\Psik(n,K)-\Psik(n,\infty)\big]\le\eta ,
\]
with the width of the shell chosen so that a photon boundary term is below $\eta$ as well. At every fixed spacing the bracket tends to $0$ as $K\to\infty$ (\Cref{rot:prop:fixedeps}); the content of \REz\ is that the heavy-shell limit and the continuum limit can be interchanged, at one geometry and with no rate. A two-line sandwich (\Cref{rot:lem:TL}) shows that the heavy shell is bounded below by a higher wall, and once the matter on the shell is dead the photon on the far side can be pinned and then replaced by King's Neumann photon on a larger cube at a cost exponentially small in the width of the shell, uniformly in the spacing (\Cref{rot:lem:GB}); so \REz\ implies \TIz\ (\Cref{rot:thm:REStoTIGHT}); in the exterior-shell geometry the same chain is \DSh\ (\Cref{rot:hyp:DS}). Under \Kz\ the limit $\lim_k\Psik(n,K)$ exists for every finite $K$, and \REz\ is equivalent to the statement that the monotone limit $\lim_K\lim_k\Psik(n,K)$ lies within $\eta$ of $\limsup_k\Psik(n,\infty)$ (\Cref{rot:thm:cont}); when the dead-shell limit $\lim_k\Psik(n,\infty)$ exists, which is a further existence statement not needed for \REz, this is the continuum form \REcont, in which no lattice spacing appears. In the curvature sector \REz\ is proved at one configuration (\Cref{rot:thm:shell}): the shell mass $\varkappa$ is put inside the matter covariance of King's block-spin flow, and two estimates on that flow, a two-sided bound on the stopping-scale integral uniformly in $\varkappa$ and a $\varkappa$-Cauchy property of the reduced small-field action on one test set, and a disc bound for the heavy shell, which is established there, give the tail bound at one configuration. The two estimates form \HS\ (\Cref{rot:hyp:HS}); they come from a renormalisation-group analysis of King's walled torus model with the shell mass inside the matter covariance (\Cref{rot:prop:SW1,rot:prop:SW3}).

The torus comparison is reduced to one-sided bounds at a single lattice spacing (\Cref{rot:thm:sandwich}). The passive map $S=R^{-1}\sigma_\nu$ is an involution that exchanges the two lattices and maps the torus $\DJ$ to the torus $R^{-1}\DJ$. A forward bound of King's type on each of the two tori therefore gives both directions of the comparison on $\DJ$. \Cref{rot:thm:sandwich} reduces \Kz\ in the curvature sector to the one-sided bounds \hypFz\ and \hypV\ of \Cref{rot:hyp:F} (bounds at one spacing, with a condition on their constants) together with the existence of the axis-lattice limits (\Cref{rot:rem:closing}); these are what the proof of \Kz\ along King's lines produces, with the large-field clause of King's upper bound in place of King's torus bound~\cite[(4.2)]{K2}, which is vacuous at a non-zero wall (\Cref{rot:rem:closing}, \Cref{rot:rem:behindK}). An alternative placement of the comparison, on a Neumann--Dirichlet box aligned with one lattice and a staircase for the other, is possible in principle and is not used here.

\subsection{The main theorem of the rotation part}

For $R\in\Oth$ and $b\in\R^3$, sources are transported as $2$-forms, $R\cdot f$ and $\tau_bf$ (\Cref{def:s02:sources}).

\begin{theorem}[$O(3)$ invariance of the curvature sector]\label{rot:thm:main}
King's printed state on curvature sources is $O(3)$-invariant: for every $R\in\Oth$, all real $2$-forms $f_1,\dots,f_N\in C^\infty_c(\R^3;\Lambda^2)$ and all $z\in\C^N$,
\[
\ZKinf\Big(\sum_jz_j\,R\cdot f_j,\,0\Big)=\ZKinf\Big(\sum_jz_jf_j,\,0\Big),
\]
so all curvature Schwinger functions are $O(3)$-covariant. With the translations of \Cref{cor:s06:TK} the same holds for every Euclidean motion $x\mapsto Rx+b$, with $\tau_bR\cdot f_j$ in place of $R\cdot f_j$: the curvature sector is invariant under $\Euc$, and the curvature Schwinger functions are Euclidean covariant.
\end{theorem}

The proof is assembled at the end of the heavy-shell sections (after \Cref{rot:thm:shell}): the heavy-shell estimates \HS\ (\Cref{rot:hyp:HS}) give \REz\ at one configuration (\Cref{rot:thm:shell}), \REz\ gives the closing estimate \TIz\ (\Cref{rot:thm:REStoTIGHT}, with the torus comparison \Kz\ of \Cref{rot:hyp:K} and the dead-shell chain \DSh\ of \Cref{rot:hyp:DS}), and \Cref{rot:thm:curv} gives the invariance from \EC\ (\Cref{hyp:s03:EC}), \Kz\ and \TIz. For the curvature two-point function the weaker form \TItwo\ of the closing estimate suffices (\Cref{rot:cor:twopoint}). The rotation part is organised as follows: this section sets up the lattices with a frame, the walled tori and the three lemmas (cutting, passive covariance, second moment) that every later argument uses; the next section treats Wick schemes and the full object on the comparison convention, which the scalar-source forms of the theorems need; then come the torus comparison, the closing hypothesis with the rotation theorems, the two heavy-shell sections, and the corrections to~\cite{K1,K2} that the rotation part relies on.

\subsection{Geometry}

\begin{definition}\label{rot:def:geom}
Let $\thz\in(0,\pi/2)$ with $\cos\thz=\frac45$, $\sin\thz=\frac35$. For $1\le\mu<\nu\le3$ let $R_{\mu\nu}(\theta)$ be the rotation with $Re_\mu=\cos\theta\,e_\mu+\sin\theta\,e_\nu$, $Re_\nu=-\sin\theta\,e_\mu+\cos\theta\,e_\nu$, $Re_\rho=e_\rho$ for $\rho\ne\mu,\nu$. Put $\Rgen\defeq\{R_{\mu\nu}(\pm\thz):1\le\mu<\nu\le3\}$ and let $\Gz$ be the group generated by $\Rgen$. The spacings $\eps_k=5^{-k}$, the axis lattices $\Lk=\eps_k\Z^3$, King's large-block integer $M\ge1$ \cite[\S2.3]{K1} and King's cubes $\CJ=[-MJ,MJ)^3$, $J\in\N$ \cite[(2.7)]{K2}, are those of the first sections. The \emph{King tori} of the rotation part are $\DJ\defeq\R^3/(50MJ\,\Z)^3$ with fundamental cell $\GamJ\defeq[-25MJ,25MJ)^3$ \cite[(4.1)]{K2}; $\DJ$ is a King torus in the sense of the section on the model ($\ell=50MJ\in M\N$ is even), so $\Lk$ descends to it. We write $\Qb b\defeq[-b,b]^3$ (written $Q_R$ in the first sections), $\Nr rX\defeq\{y:\dist_\infty(y,X)\le r\}$, and $\Bbar\rho$ for the closed Euclidean ball of radius $\rho$ about $0$.
\end{definition}

\begin{lemma}\label{rot:lem:geom}
\begin{enumerate}[label=\textup{(\alph*)},leftmargin=*]
\item For $R\in\Rgen$ the matrix $5R$ has integer entries, $R^{-1}\in\Rgen$, and $\sigma_\nu R_{\mu\nu}(\theta)\sigma_\nu=R_{\mu\nu}(-\theta)$, where $\sigma_\nu$ is the reflection $x_\nu\mapsto-x_\nu$.
\item For $R\in\Rgen$, $k\ge0$ and $J\ge1$ the lattices $R\Lk$ and $R^{-1}\Lk$ are invariant under translations by $50MJ\,\Z^3$, so they descend to $\DJ$.
\item For $R\in\Rgen$, $|Rx|_\infty\le\frac75|x|_\infty$, and $\Qb{5b/7}\subset R^{-1}\Qb b\subset\Qb{7b/5}$ for $b>0$.
\item $\thz/\pi$ is irrational, and $\Gz$ is dense in $SO(3)$.
\end{enumerate}
\end{lemma}

\begin{proof}
(a) The entries of $R$ are $0,1,\pm\frac35,\pm\frac45$, and $R_{\mu\nu}(\theta)^{-1}=R_{\mu\nu}(-\theta)$; conjugation by $\sigma_\nu$ reverses the orientation of the $(x_\mu,x_\nu)$-plane and fixes the third axis.
(b) For $p\in50MJ\,\Z^3$ we have $R^{\pm1}p=\frac15(5R^{\pm1})p\in10MJ\,\Z^3\subset\Lk$. Hence $R^{-1}\Lk+p=R^{-1}(\Lk+Rp)=R^{-1}\Lk$, and likewise for $R\Lk$.
(c) The sums of the absolute values of the entries in each row of $R$ are at most $\frac75$. If $Rx\in\Qb b$ then $|x|_\infty=|R^{-1}(Rx)|_\infty\le\frac75b$, since $R^{-1}\in\Rgen$; if $|x|_\infty\le\frac57b$ then $|Rx|_\infty\le b$.
(d) If $\thz\in\pi\mathbb Q$, then $e^{\mathrm i\thz}=(4+3\mathrm i)/5$ is a root of unity, hence an algebraic integer. Its minimal polynomial over $\mathbb Q$ is $x^2-\frac85x+1$, which does not have integer coefficients; this is a contradiction (compare~\cite[pp.~337--338]{K2}). Hence $\{n\thz\}_{n\in\Z}$ is dense modulo $2\pi$, and the closure of $\{R_{\mu\nu}(\thz)^n\}_{n\in\Z}$ is the group of all rotations in the $(x_\mu,x_\nu)$-plane. Every element of $SO(3)$ is a product $R_{12}(\alpha)R_{23}(\beta)R_{12}(\gamma)$ (Euler angles), so $\overline{\Gz}=SO(3)$.
\end{proof}

\subsection{Lattices with a frame}

\begin{definition}[lattices with a frame, plaquettes, sources]\label{rot:def:frames}
A \emph{lattice of spacing $\eps_k$} is $\Lambda=O\Lk$ with $O\in\Oth$, equipped with the frame $f_\alpha\defeq Oe_\alpha$, $\alpha=1,2,3$, and with the nearest-neighbour bonds $\langle x,x+\eps_kf_\alpha\rangle$. (The frame vectors $f_\alpha$ and the $2$-form sources $f$ share a letter; the index tells them apart.) A field configuration on a set $X$ of sites consists of $\phi(x)=r_xe^{\mathrm i\vartheta_x}\in\C$ and $A_\alpha(x)\in\R$, $x\in X$; the bond $\langle x,x+\eps f_\alpha\rangle$ carries $A_{\langle x,x+\eps f_\alpha\rangle}\defeq A_\alpha(x)$ and the reversed bond carries $-A_\alpha(x)$ (Feynman form of the vector field,~\cite[(2.5)]{K2}). The action on $\Lambda$ is King's action of the section on the model written in the frame, $e_\alpha$ replaced by $f_\alpha$ throughout, with the same free-boundary counterterm $\dm_\eps$ and the same free Wick constant $\sW_\eps$, which depend on $\eps$ only. A plaquette is $p=(x;\alpha,\beta)$ with $\alpha<\beta$; its centre is $c(p)\defeq x+\frac\eps2(f_\alpha+f_\beta)$ and its field strength is
\[
F(p)\defeq\eps^{-1}\big(A_\alpha(x)+A_\beta(x+\eps f_\alpha)-A_\alpha(x+\eps f_\beta)-A_\beta(x)\big)
\qquad\text{\cite[(2.11)]{K2}}.
\]
For a real $2$-form $f$ and a real function $h$, smooth and compactly supported (on a torus: smooth and periodic; a compactly supported source whose support lies in the interior of a fundamental cell is periodised from it), put
\[
F_\Lambda(f)\defeq\sum_p\eps^3\,f_{c(p)}(f_\alpha,f_\beta)\,F(p),\qquad
\sS_\Lambda(h)\defeq\sum_x\eps^3h(x),\qquad
W_\Lambda(h)\defeq\sum_x\eps^3h(x)\big(|\phi(x)|^2-2\sW_\eps\big),
\]
the sums running over the plaquettes $p=(x;\alpha,\beta)$, respectively the sites, of the region (of one fundamental cell on a torus): the $2$-form is \emph{sampled in the lattice's own frame}. On $\Lk$ with the standard frame these are $F_{\Lk}(f)$, $\sS_{\Lk}(h)$ and $\Wick{h}$ of \Cref{def:s02:sources}. For a finite family $\zeta=(z_j,f_j;w_j,h_j)_j$ with complex coefficients, $\Theta(\zeta)\defeq\exp(\sum_jz_jF_\Lambda(f_j)+\sum_jw_jW_\Lambda(h_j))$ as there, and $O\cdot\zeta\defeq(z_j,O\cdot f_j;w_j,O\cdot h_j)_j$ with the transport of \Cref{def:s02:sources}. The cone $\cK$ of \Cref{def:s02:sources} is invariant under $\Oth$, under translations, under $(f,g)\mapsto(tf,tg)$, $t\ge0$, and under sums.
\end{definition}

\begin{remark}[the frame convention]\label{rot:rem:frame}
In~\cite{K1,K2} the curvature source is written $F(g)=\sum_p\eps^3g(p)F(p)$ with one function $g$ on all plaquettes \cite[(2.12)]{K2}. On one lattice a $2$-form sampled as above is a function on plaquettes, and King's proofs treat the source terms componentwise (``the sources $g_{\mu\nu}(x'),h(x')$ \dots\ can be replaced by $g_{\mu\nu}(x),h(x)$'', \cite[p.~665]{K1}; the sentence is stated there without an argument). King's single-lattice statements at $h=0$ are used in this sense throughout: the existence theorem \EC\ is stated, in \Cref{hyp:s03:EC}, for real smooth $2$-form sources sampled at plaquette centres in the lattice's own frame, and this existence clause of King's Theorem~2.1 of~\cite{K2} is proved in \cref{sec:k01}; the statements of~\cite[Thms.~2.2, 2.3, (3.40)--(3.44)]{K2} that the rotation part uses, all at $h=0$, are established in the first sections from \EC\ and the correlation inequalities, \cite[(3.42)--(3.43)]{K2} are not used, and the torus statement~\cite[Thm.~2.1]{K1} at a wall is part of the torus comparison \Kz\ (\Cref{rot:hyp:K}). The two-lattice statement is different: \Cref{rot:prop:scalar} shows that it fails for a single scalar function on all plaquettes, and \Kz\ is formulated for $2$-forms sampled in each lattice's own frame.
\end{remark}

\subsection{Tori and walled tori}

\begin{definition}[periodic and walled torus states]\label{rot:def:tori}
Let $\Lambda$ be a lattice of spacing $\eps=\eps_k$ with a frame, and $d\nu\defeq\prod_x\prod_\alpha dA_\alpha(x)\,d^2\phi(x)$ over the sites in question.
\begin{enumerate}[label=\textup{(\alph*)},leftmargin=*]
\item \emph{Boxes.} For a finite set $X$ of sites of $\Lambda$, $\langle\cdot\rangle^{\ND}_{\Lambda,X}$ is the Neumann--Dirichlet state of \Cref{def:s02:acube} written in the frame of $\Lambda$ (Dirichlet condition for the scalar field outside $X$, Neumann condition for the vector field), possibly with the matter killed on a set $U\subset X$ in the sense of that definition. On $\Lk$ these are the states $\langle\cdot\rangle^{\ND}_{k,X}$ of \Cref{def:s02:acube}, among them the comparison cube $\ava\cdot$ and King's RG cube $\avK\cdot$ (\Cref{def:s02:Kcube}).
\item \emph{Tori.} If $\Lambda$ descends to $\DJ$, $\langle\cdot\rangle_{\Lambda,\DJ}$ is the normalised measure $Z^{-1}e^{-S}d\nu$ with the action $S$ summed over the sites and bonds of the finite graph $\Lambda/(50MJ\,\Z)^3$ (periodic conditions for both fields). On $\Lk$ this is the periodic state $\langle\cdot\rangle_{k,T}$ of the section on the model for the King torus $T=\DJ$.
\item \emph{Walled tori.} A \emph{walled configuration} is $\fc=(J',w)$ with $J'\in\N$ and $w\in C^\infty(D_{J'};[0,\infty))$, the \emph{wall}. If $\Lambda$ descends to $D_{J'}$, $\langle\cdot\rangle_{\fc,\Lambda}$ is the torus state of (b) with the density multiplied by $\prod_x\exp(-\eps^3w(x)|\phi(x)|^2)$. The wall raises the scalar mass squared by $2w$ and acts in the single-site measures. We write $(J',\psi,n)\defeq(J',n\psi)$ for $n\ge0$ and a fixed profile $\psi$.
\item \emph{Killed sites.} For any of the states (b), (c) and a set $U$ of sites, the state \emph{with the matter killed on $U$} is defined as in \Cref{def:s02:acube}: the single-site measure at $x\in U$ is replaced by the point mass at $\phi(x)=0$, every hopping term with an endpoint in $U$ reduces to its diagonal part at the other endpoint, and the photon variables are unchanged.
\end{enumerate}
For every state $\langle\cdot\rangle_{\Lambda,X}$ above and a scheme $a$ (\Cref{def:s02:sources}, \Cref{rem:s02:scheme}) the scheme functional is defined by the formula of \Cref{def:s02:sources} with $F_\Lambda$, $W_\Lambda$, $\sS_\Lambda$ in place of their axis-lattice versions,
\[
Z^a_{\Lambda,X}(\zeta)\defeq e^{-\sum_jw_ja_k\sS_\Lambda(h_j)}\big\langle\Theta(\zeta)\big\rangle_{\Lambda,X};
\]
on $\Lk$ we keep the notation $Z^a_{k,X}$ of \Cref{def:s02:sources}, and on the rotated lattice we write $Z^a_{R^{-1}\Lk,X}$. All these states are invariant under the charge conjugation $(A,\vartheta)\mapsto(-A,-\vartheta)$, which maps $F_\Lambda(f)$ to $-F_\Lambda(f)$ and fixes $r$; hence $\langle\Theta(\mathrm if,g)\rangle=\langle\cos F_\Lambda(f)\,e^{W_\Lambda(g)}\rangle$ is real for $(f,g)\in\cK$.
\end{definition}

\begin{definition}[the unnormalised torus functional and the walled ratio]\label{rot:def:Zun}
For a lattice $\Lambda$ of spacing $\eps_k$ descending to a torus $T$, and real smooth periodic sources $G$ (a $2$-form) and $h$ (a function), the \emph{unnormalised} generating functional is
\[
\Zun_{\Lambda,T}(G,h)\defeq\int e^{-S(A,\phi)+F_\Lambda(G)+W_\Lambda(h)}\,d\nu ,
\]
with no field-independent constant in the action; for a pair $s=(G,h)$ we also write $\Zun_{\Lambda,T}(s)$. Two lattices of the same spacing descending to $T$ have the same number of sites in a fundamental cell, so any field-independent constant per site multiplies both by the same factor. For a walled configuration $\fc=(J',w)$, a lattice $\Lambda$ descending to $D_{J'}$, a scheme $a$ and $(f,g)\in\cK$ with $w=0$ on the support of $g$, the \emph{walled ratio} is
\[
\Tw a{\fc,\Lambda}(f,g)\defeq e^{-a_k\sS_\Lambda(g)}\big\langle\Theta(\mathrm if,g)\big\rangle_{\fc,\Lambda}
=\frac{Z^a_{\Lambda,D_{J'}}(\mathrm if,g-w)}{Z^a_{\Lambda,D_{J'}}(0,-w)} .
\]
The second expression agrees with the first because
\[
e^{-W_\Lambda(w)}=e^{2\sW_\eps\sS_\Lambda(w)}\prod_xe^{-\eps^3w(x)|\phi(x)|^2},
\]
so that the wall is the scalar source $-w$ up to the constant $e^{2\sW_\eps\sS_\Lambda(w)}$, which cancels between numerator and denominator, as do the multipliers $e^{a_k\sS_\Lambda(w)}$ of both. The same formulas define the walled ratio for the periodic model of $\Lambda$ on any torus $T=\R^3/P$ to which $\Lambda$ descends (defined as in \Cref{rot:def:tori}(b), with the finite graph $\Lambda/P$) and any smooth wall $w\ge0$ on $T$ vanishing on the support of $g$; this is used only for the transported tori of \Cref{rot:lem:walledalg}(b).
\end{definition}

Every state of \Cref{rot:def:tori}, with or without killed sites, is a BFS measure of King type in the sense of the section on the model: $c_{xy}=\eps$ for the bonds present, the single-site measures contain $e^{-\eps^3[\frac12(m^2+\dm_\eps)r^2+\lambda r^4]}$, the diagonal parts $\frac\eps2r_x^2$ of all bonds at $x$ (including, for boxes, those leading out of $X$, and those leading into a killed site) and the wall factor $e^{-\eps^3w(x)r_x^2}$ with $w\ge0$; the photon form is positive definite because $\lmu>0$, and the exponential moments are finite because $\lambda>0$. The same holds for every measure obtained from such a state by multiplying some $c_{xy}$ by $t\in[0,1]$, multiplying photon gradient terms by $t\in[0,1]$, adding $s|\ell(A)|^2$ ($s\ge0$) to the action, multiplying single-site measures by $e^{-sr^2}$ ($s\ge0$), or killing sites. Curvature sources $F_\Lambda(f)$ are real linear functionals of $A$, and $e^{W_\Lambda(g)}$ is a positive constant times $e^{\sum_x\eps^3g(x)r_x^2}$; so \Cref{lem:s02:BFS}, \Cref{lem:s02:GD} and \Cref{lem:s02:Kprime} apply to all of them, on every lattice with a frame.

\subsection{Cutting, passive covariance and the second moment}

\begin{lemma}[cutting, walled tori]\label{rot:lem:cutwall}
Let $\langle\cdot\rangle_Y$ be a box, torus or walled-torus state of \Cref{rot:def:tori} on a lattice $\Lambda$ with a frame, possibly with killed sites, and let $X$ be a set of sites of $\Lambda$ such that $X$ and its nearest neighbours lie in one fundamental cell. Let $\langle\cdot\rangle^{\ND}_X$ be the Neumann--Dirichlet state on $X$ whose single-site measures (wall, killed sites and diagonal parts of the bonds leaving $X$ included) are those of $\langle\cdot\rangle_Y$. Let $(f,g)\in\cK$ be such that every plaquette and every site at which $f$, respectively $g$, does not vanish lies in $X$. Then, at every lattice spacing,
\[
\langle\Theta(\mathrm if,g)\rangle^{\ND}_X\le\langle\Theta(\mathrm if,g)\rangle_Y,\qquad
\langle\Theta(f,-g)\rangle^{\ND}_X\ge\langle\Theta(f,-g)\rangle_Y ,
\]
and the same for $Z^a$. If the wall of $Y$ vanishes on $X$ and no site of $X$ is killed, $\langle\cdot\rangle^{\ND}_X=\langle\cdot\rangle^{\ND}_{\Lambda,X}$ is the box state of \Cref{rot:def:tori}(a).
\end{lemma}

\begin{proof}
The wall and the killed sites are part of the single-site measures, and the frame plays no role, so the proof of \Cref{lem:s02:cut} applies word for word. Delete from $Y$ every photon gradient term coupling a site of $X$ to a site outside $X$, and the cosine parts of all hopping terms between $X$ and its complement, keeping their diagonal parts $\frac\eps2(r_x^2+r_y^2)$; since $X$ and its neighbours lie in one fundamental cell, these are bonds between distinct sites of the finite graph. By \Cref{lem:s02:BFS}(a),(b) this does not increase the first expectation, by (d) it does not decrease the second, and every intermediate measure is a BFS measure. The resulting measure is a product (\Cref{lem:s02:BFS}(e)); its $X$-factor is $\langle\cdot\rangle^{\ND}_X$, since the diagonal parts of the cut bonds are exactly the Dirichlet terms and the remaining photon terms are the Neumann ones, and the observable depends on the $X$-variables only. The free Wick constant, $\dm_\eps$ and the scheme multiplier $e^{\mp a_k\sS_\Lambda(g)}$ are the same in both states.
\end{proof}

With $Y=X_{J+1}$ and $X=\XJ$ on $\Lk$ this is King's monotonicity in the volume~\cite[Thm.~3.2, (3.31)--(3.33)]{K2} in the form his proof proves, with one Wick constant for all $J$, as discussed after \Cref{lem:s02:cut}; with $Y$ a walled torus and $X$ a cube of $\Lambda$ inside the cell, it is the inequality of~\eqref{rot:eq:chain}.

\begin{lemma}[passive covariance and the algebra of the walled ratio]\label{rot:lem:walledalg}
Let $a$ be a scheme.
\begin{enumerate}[label=\textup{(\alph*)},leftmargin=*]
\item For real coefficients the multiplier $e^{-\sum_jw_ja_k\sS_\Lambda(h_j)}$ is a positive constant depending only on $k$, the coefficients and the numbers $\sS_\Lambda(h_j)$. Hence the inequalities of \Cref{lem:s02:BFS} and \Cref{rot:lem:cutwall} between two states on the same lattice with the same sources hold for $Z^a$, and for $\Tw a{\fc,\Lambda}$, in the same direction.
\item Let $O\in\Oth$, let $\Lambda$ be a lattice of spacing $\eps$ with frame $(f_\alpha)$, and let $X$ be a box, a torus or a walled torus for $\Lambda$, possibly with the matter killed on $U$. Equip $O\Lambda$ with the frame $(Of_\alpha)$ and $OX$ with the transported boundary conditions, the wall $O\cdot w$ and the killed set $OU$; for a torus $X=\R^3/P$ this means the torus $OX=\R^3/OP$, to which $O\Lambda$ descends. Then for every finite family $\zeta$
\[
\big\langle\Theta(O\cdot\zeta)\big\rangle_{O\Lambda,OX}=\big\langle\Theta(\zeta)\big\rangle_{\Lambda,X},\qquad
\sS_{O\Lambda}(O\cdot h)=\sS_\Lambda(h),\qquad
Z^a_{O\Lambda,OX}(O\cdot\zeta)=Z^a_{\Lambda,X}(\zeta).
\]
In particular, for a walled configuration $\fc=(J',w)$ and $(f,g)$ as in \Cref{rot:def:Zun}, $\Tw a{O\cdot\fc,O\Lambda}(O\cdot f,O\cdot g)=\Tw a{\fc,\Lambda}(f,g)$, where $O\cdot\fc$ denotes the transported walled torus, the torus $OD_{J'}$ with the wall $O\cdot w$, and the left side is the walled ratio on that torus in the sense of \Cref{rot:def:Zun}; $O\cdot\fc$ is a walled configuration $(J',O\cdot w)$ of \Cref{rot:def:tori}(c) when $OD_{J'}=D_{J'}$, as for a signed permutation matrix $O$. Moreover, on every torus the periodic model on $\Lambda$ does not change if the frame $(f_\alpha)$ is replaced by $(\pm f_{\pi(\alpha)})$, for a permutation $\pi$ and any choice of signs.
\item $\Tw a{\fc,\Lambda}(f,g)=e^{-a_k\sS_\Lambda(g)}\langle\Theta(\mathrm if,g)\rangle_{\fc,\Lambda}$ is the normalised expectation of the walled torus state up to the multiplier of \textup{(a)}, and at $g=0$ the walled ratio $\Tw a{\fc,\Lambda}(f,0)=\langle e^{\mathrm iF_\Lambda(f)}\rangle_{\fc,\Lambda}$ does not depend on the scheme.
\end{enumerate}
\end{lemma}

\begin{proof}
(a) The first statement is immediate (compare the remark in~\cite[p.~281]{BFS1} that the inequalities are not affected by Wick ordering); for $\Tw a{\fc,\Lambda}$ the multiplier is $e^{-a_k\sS_\Lambda(g)}$.
(b) The map $\phi'(Ox)=\phi(x)$, $A'_\alpha(Ox)=A_\alpha(x)$, with components taken in the frames $(f_\alpha)$ and $(Of_\alpha)$, is a bijection of the configuration spaces preserving $d\nu$; this is the proof of \Cref{lem:s02:pass}, which covers boxes and tori, and for a torus $X=\R^3/P$ it carries the finite graph $\Lambda/P$ onto $O\Lambda/OP$. The action is written in the frame, and $\dm_\eps$, $\sW_\eps$ and $a_k$ depend on $\eps$ only; the wall term $\sum_x\eps^3w(x)|\phi(x)|^2$ becomes $\sum_{x}\eps^3(O\cdot w)(Ox)|\phi'(Ox)|^2$ since $(O\cdot w)(Ox)=w(x)$, and a point mass at $\phi(x)=0$ becomes a point mass at $\phi'(Ox)=0$. Plaquettes, their centres and field strengths are transported, and $(O\cdot f)_{Oc}(Of_\alpha,Of_\beta)=f_c(f_\alpha,f_\beta)$, $\sum_{Ox}\eps^3(O\cdot h)(Ox)=\sum_x\eps^3h(x)$. The clause on the walled ratio is the first identity for $\zeta=(\mathrm i,f;1,g)$ and the walled torus $X=D_{J'}$, together with $\sS_{O\Lambda}(O\cdot g)=\sS_\Lambda(g)$ and the first expression of \Cref{rot:def:Zun} read on $OD_{J'}$. For the last statement, a permutation of the frame relabels the components; reversing $f_\alpha$ replaces $A_\alpha(x)$ by $-A_\alpha(x-\eps f_\alpha)$, a bijection of the torus variables that preserves the action, and changes the sign of both $F(p)$ and $f_{c(p)}(f_\alpha,f_\beta)$ while preserving the set of plaquette centres.
(c) The expression for $\Tw a{\fc,\Lambda}$ is its definition, and at $g=0$ the multiplier is $1$ and $W_\Lambda(0)=0$.
\end{proof}

\begin{lemma}[the characteristic function and the second moment of a photon functional]\label{rot:lem:GDb}
Let $\langle\cdot\rangle$ be a BFS measure whose Gaussian has covariance $C$, and let $\Psi(A)=\sum_qw(q)\cos\ip{\kappa(q)}A$ and $\E^\sharp$ be the current representation and the current measure of \Cref{lem:s02:GD}. Then for $v\in\R^m$:
\begin{enumerate}[label=\textup{(\alph*)},leftmargin=*]
\item $\langle e^{\mathrm i\ip vA}\rangle=e^{-\ip v{Cv}/2}\,\E^\sharp\cosh\ip{Cv}{\kappa}$; in particular $\langle e^{\mathrm i\ip vA}\rangle\in[0,1]$ is real;
\item $\langle\ip vA\rangle=0$ and $\langle\ip vA^2\rangle=\ip v{Cv}-\E^\sharp\ip{Cv}{\kappa}^2\in[0,\ip v{Cv}]$.
\end{enumerate}
\end{lemma}

\begin{proof}
(a) For a centred Gaussian $\gamma$ with covariance $C$ and $a\in\C^m$, $\int e^{\ip aA}d\gamma=e^{\ip a{Ca}/2}$. With $a=\mathrm i(v\pm\kappa)$ and the average of the two signs,
\[
\int e^{\mathrm i\ip vA}\cos\ip\kappa A\,d\gamma=e^{-\ip v{Cv}/2-\ip\kappa{C\kappa}/2}\cosh\ip v{C\kappa},
\]
since $C$ is symmetric. Multiply by $w(q)\ge0$, sum over the currents $q$ (the series converges absolutely: $\cosh\ip v{C\kappa}\le e^{|\ip v{C\kappa}|}\le e^{\ip v{Cv}/2+\ip\kappa{C\kappa}/2}$ by the Cauchy--Schwarz inequality for the form $C$, so the terms are dominated by $w(q)$ and $\sum_qw(q)<\infty$), and divide by the same sum at $v=0$, which is the positive normalisation $\sum_qw(q)e^{-\ip\kappa{C\kappa}/2}$ of $\E^\sharp$; this is the formula. Its right side lies in $[0,\infty)$, and $|\langle e^{\mathrm i\ip vA}\rangle|\le1$.
(b) Replace $v$ by $tv$, $t\in\R$, in (a). The left side is the characteristic function of the real random variable $\ip vA$, which has all moments because $\langle e^{s\ip vA}\rangle\le e^{s^2\ip v{Cv}/2}$ for every real $s$ (\Cref{lem:s02:GD}); so it equals $1+\mathrm it\langle\ip vA\rangle-\frac{t^2}2\langle\ip vA^2\rangle+O(t^3)$. The right side is $e^{-t^2\ip v{Cv}/2}\E^\sharp\cosh(t\ip{Cv}\kappa)$, and $\E^\sharp\cosh(t\ip{Cv}\kappa)=1+\frac{t^2}2\E^\sharp\ip{Cv}\kappa^2+O(t^4)$, the expansion under $\E^\sharp$ being justified by the domination $\cosh(t\ip{Cv}\kappa)\le e^{t^2\ip v{Cv}/2+\ip\kappa{C\kappa}/2}$ of (a), which shows that $\ip{Cv}\kappa$ has all moments under $\E^\sharp$. The right side is real and even in $t$, so $\langle\ip vA\rangle=0$; comparing the terms of order $t^2$ gives $\langle\ip vA^2\rangle=\ip v{Cv}-\E^\sharp\ip{Cv}\kappa^2$. The bounds follow since both $\langle\ip vA^2\rangle$ and $\E^\sharp\ip{Cv}\kappa^2$ are non-negative.
\end{proof}

The lemma applies to every state of \Cref{rot:def:tori} and to every state of the first sections, with $v$ the vector $\dstar_\Lambda f$ defined by $F_\Lambda(f)=\ip{\dstar_\Lambda f}A$ for a real $2$-form $f$ whose plaquettes carry variables of the state; the photon Gaussian of a torus or walled torus is the periodic Feynman form with mass $\lmu$, and killed sites are covered as in \Cref{lem:s02:GD}.
